\documentclass[12pt,a4paper,UTF8]{ctexart}
\usepackage[a4paper,top=22mm,bottom=22mm,left=22mm,right=22mm,headheight=15pt]{geometry}
\newcommand{\Needspace}[1]{\par\begingroup\dimen0=#1\relax\dimen2=\pagegoal\advance\dimen2 by-\pagetotal\ifdim\dimen0>\dimen2\newpage\fi\endgroup}
\newcommand{\needspace}[1]{\par\begingroup\skip0=\lastskip\vskip-\skip0\dimen0=#1\relax\advance\dimen0 by\skip0\vskip\dimen0\penalty-100\vskip-\dimen0\vskip\skip0\penalty9999\endgroup}
\usepackage{amsmath,amssymb,mathtools,bm}
\usepackage{graphicx,xcolor,booktabs,longtable,array,enumitem,fancyhdr,caption,float,microtype,setspace}
\usepackage[hidelinks,unicode]{hyperref}
\usepackage[most]{tcolorbox}
\definecolor{ContestBlue}{HTML}{17365D}
\hypersetup{pdftitle={第054套},pdfauthor={}}
\setstretch{1.24}
\setlength{\parindent}{2em}
\setlength{\parskip}{0.35em}
\setlength{\emergencystretch}{2em}
\setlist{itemsep=0.35em,topsep=0.4em}
\captionsetup{font=small,labelfont=bf,justification=centering}
\setcounter{secnumdepth}{0}
\pagestyle{fancy}
\fancyhf{}
\fancyhead[L]{\small 第054套}
\fancyfoot[C]{\thepage}
\renewcommand{\headrulewidth}{0.4pt}
\ctexset{
  section={format=\Large\bfseries\color{ContestBlue},beforeskip=1.2em,afterskip=0.7em},
  subsection={format=\large\bfseries\color{ContestBlue},beforeskip=1.0em,afterskip=0.5em},
  subsubsection={format=\normalsize\bfseries,beforeskip=0.8em,afterskip=0.35em}
}
\usepackage{calc}
\newcounter{none}
\let\originalsubsection\subsection
\renewcommand{\subsection}{\needspace{6\baselineskip}\originalsubsection}
\makeatletter
\newcommand{\tightlist}{\setlength{\itemsep}{0pt}\setlength{\parskip}{0pt}}
\newcommand{\maxwidth}{\ifdim\Gin@nat@width>\linewidth\linewidth\else\Gin@nat@width\fi}
\newcommand{\maxheight}{\ifdim\Gin@nat@height>0.72\textheight 0.72\textheight\else\Gin@nat@height\fi}
\newcommand{\pandocbounded}[1]{\begingroup\setkeys{Gin}{width=\maxwidth,height=\maxheight,keepaspectratio}#1\endgroup}
\makeatother
\begin{document}
\begin{titlepage}
  \centering
  \vspace*{58mm}
  {\Huge\bfseries\color{ContestBlue} 第054套\par}
  \vfill
  \vspace*{28mm}
\end{titlepage}
\section{第 1 题}\label{ux7b2c-1-ux9898}

水平面内有一根质量均匀分布的光滑细丝，形状为等角螺线

\[
r=r_0e^{\theta/\sqrt3},\qquad 0\le\theta\le2\pi,
\]

其中\(r_0=1\ \mathrm m\)，总质量为
\(m=1\ \mathrm{kg}\)。细丝由轻质支架固定，可绕过原点且垂直于平面的轴无摩擦转动。

（1）求螺线总长度 \(l\)；

（2）求螺线对转轴的转动惯量 \(I\)；

（3）一个质量同为 \(m\) 的小物体套在 \(\theta=0\)
端，沿螺线切向获得初速度
\(v_0=10\ \mathrm{m/s}\)，螺线初始静止。求小物体从另一端离开时螺线的角速度；

（4）若空间中再有垂直纸面向内的匀强磁场 \(B=1\ \mathrm T\)，小物体带正电
\(q=0.01\ \mathrm C\)，螺线绝缘，重新求第（3）问的角速度。

\begin{figure}[H]
\centering
\includegraphics[width=0.8\linewidth,height=\textheight,keepaspectratio,alt={图1-1　等角螺线、端点半径与初速度方向}]{../figures/question-01.pdf}
\caption*{图1-1　等角螺线、端点半径与初速度方向}
\end{figure}

\newpage

\section{第 2 题}\label{ux7b2c-2-ux9898}

一列简化火车由两个半径为 \(r\)、质量均为 \(m\) 的实心轮和一根长度为
\(3r\)、质量为 \(3m\) 的均匀车身组成。断桥宽度为
\(L=10r\)，桥面不提供支持力。火车初速度

\[
v=12.557\sqrt{gr},
\]

初态两轮纯滚动、无动力。车轮与右侧棱角的碰撞完全弹性，滑动摩擦系数和最大静摩擦系数均为
\(\mu=0.05\)。判断后轮能否接触右崖上部。

\begin{figure}[H]
\centering
\includegraphics[width=0.86\linewidth,height=\textheight,keepaspectratio,alt={图2-1　双轮火车飞越宽为 10r 的断桥}]{../figures/question-02.pdf}
\caption*{图2-1　双轮火车飞越宽为 \(10r\) 的断桥}
\end{figure}

\newpage

\section{第 3 题}\label{ux7b2c-3-ux9898}

\(\nu\ \mathrm{mol}\) 理想气体的摩尔等容热容为
\(C_V=5R/2\)，经历准静态循环：\(AB\) 满足 \(p=\alpha V\)，\(BC\)
绝热，\(CA\) 等温。状态 \(A,B\) 的温度分别为 \(T_1,T_2\)。求：

（1）\(AB\) 过程摩尔热容 \(C_1/R\)；

（2）\(AB\) 的吸热量 \(Q_1\)；

（3）\(CA\) 的放热量 \(Q_2\)；

（4）循环效率 \(\eta\)。

\begin{figure}[H]
\centering
\includegraphics[width=0.69\linewidth,height=\textheight,keepaspectratio,alt={图3-1　直线、绝热与等温过程组成的循环}]{../figures/question-03.pdf}
\caption*{图3-1　直线、绝热与等温过程组成的循环}
\end{figure}

\newpage

\section{第 4 题}\label{ux7b2c-4-ux9898}

边长为 \(L\) 的正方形内绘制伪希尔伯特曲线 \(H_n\)。\(H_1\)
依次连接四个小正方形中心；每升一阶，把前一阶曲线缩小并按规定旋转、平移成四份，再以三条等长线段连接。每阶曲线均由均匀细线组成，总质量均为
\(M\)。求：

（1）\(H_1\) 对 \(x\) 轴的转动惯量；

（2）\(H_1\) 对 \(y\) 轴的转动惯量；

（3）\(H_2\) 对 \(x\) 轴的转动惯量；

（4）\(H_2\) 对 \(y\) 轴的转动惯量；

（5）\(H_\infty\) 对 \(x\) 轴的转动惯量；

（6）\(H_\infty\) 对 \(y\) 轴的转动惯量。

\begin{figure}[H]
\centering
\includegraphics[width=0.88\linewidth,height=\textheight,keepaspectratio,alt={图4-1　一阶、二阶伪希尔伯特曲线与坐标轴}]{../figures/question-04.pdf}
\caption*{图4-1　一阶、二阶伪希尔伯特曲线与坐标轴}
\end{figure}

\newpage

\section{第 5 题}\label{ux7b2c-5-ux9898}

图示理想 LC 网络满足 \(C_1\ll C_2\)、\(L_1\ll L_2\)。复阻抗写作

\[
\widetilde Z=Z\,e^{i\varphi},\qquad-\pi<\varphi<\pi.
\]

求高频条件 \(\omega_1L_1\sim1/(\omega_1C_1)\) 下端口 \(DE\)
的阻抗振幅、相位和谐振频率；再求低频条件
\(\omega_2L_2\sim1/(\omega_2C_2)\) 下的对应三项。

\begin{figure}[H]
\centering
\includegraphics[width=0.431\linewidth,height=\textheight,keepaspectratio,alt={图5-1　由两组电感、电容构成的对称桥式网络}]{../figures/question-05.pdf}
\caption*{图5-1　由两组电感、电容构成的对称桥式网络}
\end{figure}

\newpage

\section{第 6 题}\label{ux7b2c-6-ux9898}

在玻尔模型中，电子质量为 \(m\)、电荷量大小为 \(e\)，静电常量为
\(k\)，圆轨道角动量为
\(n\hbar\)。索莫菲尔德模型允许以原子核为焦点的椭圆轨道，能量仍为
\(E_n\)，角动量为 \(l\hbar\)。求玻尔半径 \(r_n\)、速度 \(v_n\)、能量
\(E_n\)，椭圆最近、最远距离，偏心率，并估计相对论引起的精细结构能量差数量级。

\begin{figure}[H]
\centering
\includegraphics[width=0.606\linewidth,height=\textheight,keepaspectratio,alt={图6-1　库仑场中的圆轨道与椭圆轨道}]{../figures/question-06.pdf}
\caption*{图6-1　库仑场中的圆轨道与椭圆轨道}
\end{figure}

\newpage

\section{第 7 题}\label{ux7b2c-7-ux9898}

一理想凸透镜焦距为 \(f\)。边长为 \(f/2\) 的正方形 \(ABCD\)
位于主光轴所在平面内，顶点 \(D\) 在主轴上且距透镜
\(2f\)，正方形相对主轴倾斜 \(30^\circ\)。物方坐标原点在左焦点，\(x\)
轴沿主轴向左、\(y\) 轴向上；像方坐标原点在右焦点，\(x'\) 轴向右、\(y'\)
轴向下。求像的形状、四顶点坐标和面积。

\begin{figure}[H]
\centering
\includegraphics[width=0.879\linewidth,height=\textheight,keepaspectratio,alt={图7-1　倾斜正方形、凸透镜与焦点坐标系}]{../figures/question-07.pdf}
\caption*{图7-1　倾斜正方形、凸透镜与焦点坐标系}
\end{figure}

\newpage

\section{第 8 题}\label{ux7b2c-8-ux9898}

两根无限长直导线位于 \(yz\) 平面内，相距
\(2a\)，通有大小相等、方向相反的恒定电流 \(I\)。中央有一个尺寸为
\(2b\times d\)（\(b<a\)）的均匀矩形线框，线框电流也恒为
\(I\)，可绕通过中心的 \(z\) 轴自由转动。线框平面与 \(yz\) 平面的夹角为
\(\phi\)。求：

（1）绕 \(z\) 轴力矩大小；

（2）\(0\le\phi\le\pi/2\) 内力矩最大时的 \(\phi\)；

（3）稳定平衡角；

（4）在稳定平衡附近的小振动周期。线框质量线密度为 \(\lambda\)。

\begin{figure}[H]
\centering
\includegraphics[width=0.599\linewidth,height=\textheight,keepaspectratio,alt={图8-1　反向长直电流与可绕轴转动的矩形线框}]{../figures/question-08.pdf}
\caption*{图8-1　反向长直电流与可绕轴转动的矩形线框}
\end{figure}
\end{document}
